\section{Összegzés}

A feladat célja egy olyan mobilalkalmazás létrehozása volt, amely Parkinson-kóros
betegek rehabilitációs munkameneteinek informatikai támogatását valósítja meg.
A konfigurációs képernyőn mind az előre rögzített,
mind a teljesen egyedi rehabilitációs feladatsorok összeállíthatók. A végrehajtási képernyőn a
beépített videólejátszó és a kamera-előnézet osztott képernyős integrációjával a betegek vizuális mintát
kapnak a mozdulatok helyes kivitelezéséhez.

A szenzorból valós időben nyerhetők ki az adatok; a rendszer stabil BLE-kapcsolatot épít fel az
okoszoknival, és a kapott információkat strukturált formátumban tárolja el.

A nyers adatok feldolgozásra kerülnek, majd a páciens számára is érthető módon jelennek meg,
így a gyakorlatok végrehajtásának minősége és a rehabilitációs folyamat objektíven nyomon
követhetővé válik. Az adatok metődnek, így a korábbi munkamenet visszakövethető.

A szoftverarchitektúrában az egyes rétegek jól elkülönülnek, amellyel megvalósult a moduláris felépítés
célja. Az alkalmazás ezáltal rugalmasan karbantartható marad, és az előző fejezetben megfogalmazott
jövőbeli funkciókkal könnyedén bővíthető anélkül, hogy a meglévő logikát jelentősen módosítani kellene.

Összességében az alkalmazás a gyakorlatban is használható eszközzé vált,
amely jelenleg a tesztelési fázisban tart. Később felhasználói visszajelzések alapján
tovább finomhangolható.